% from ega4-16-fr.tex
\oldpage[IV-4]{37}
$\Delta_f:X\to X\times_S X$, par
$h_n:X_{\Delta_f}^{(n)}\to X\times_S X$ le morphisme canonique (16.1.2), et
considérons les deux morphismes composés
\[
  p_1^{(n)}:X_{\Delta_f}^{(n)}\xrightarrow{h_n}X\times_S X
  \xrightarrow{p_1}X,qquad
  p_2^{(n)}:X_{\Delta_f}^{(n)}\xrightarrow{h_n}X\times_S X
  \xrightarrow{p_2}X
\]
de sorte que, par définition, $p_1^{(n)}$ correspond à l'homomorphisme de
faisceaux d'anneaux $\mathcal{O}_X\to\mathscr{P}_{X/S}^n$ que nous avons
choisi pour définir la structure de $\mathcal{O}_X$-Algèbre sur
$\mathscr{P}_{X/S}^n$ (16.3.5), et $p_2^{(n)}$ à l'homomorphisme de faisceaux
d'anneaux $d_{X/S}^n:\mathcal{O}_X\to\mathscr{P}_{X/S}^n$ (16.3.6). Comme
$X_{\Delta_f}^{(n)}$ et $X$ ont même espace sous-jacent, on peut écrire
\[
  \mathscr{P}_{X/S}^n
  =(p_1^{(n)})_*((p_2^{(n)})^*(\mathcal{O}_X)).
  \tag{16.7.1.1}
  \label{IV.16.7.1.1}
\]
Plus généralement, nous poserons
\[
  \mathscr{P}_{X/S}^n(\mathscr{F})
  =(p_1^{(n)})_*((p_2^{(n)})^*(\mathscr{F}))
  \tag{16.7.1.2}
  \label{IV.16.7.1.2}
\]
de sorte que
$\mathscr{P}_{X/S}^n=\mathscr{P}_{X/S}^n(\mathcal{O}_X)$; par définition,
$\mathscr{P}_{X/S}^n(\mathscr{F})$ est donc un $\mathcal{O}_X$-Module.
\end{env}

\begin{env}[16.7.2]
\label{IV.16.7.2}
Si l'on revient aux définitions des images réciproques de Modules sur les
espaces annelés (0, 4.3.1) et que l'on tienne compte de ce que
$X_{\Delta_f}^{(n)}$ et $X$ ont même espace sous-jacent, on voit que l'on peut
aussi écrire la définition (16.7.1.2) sous la forme
\[
  \mathscr{P}_{X/S}^n(\mathscr{F})
  =\mathscr{P}_{X/S}^n\otimes_{\mathcal{O}_X}\mathscr{F}
  \tag{16.7.2.1}
  \label{IV.16.7.2.1}
\]
mais où il faut prendre garde que, dans l'interprétation du signe $\otimes$,
$\mathscr{P}_{X/S}^n$ est muni de sa structure de $\mathcal{O}_X$-Module
définie par l'homomorphisme de faisceaux d'anneaux
$d_{X/S}^n:\mathcal{O}_X\to\mathscr{P}_{X/S}^n$. Il résulte aussitôt de cette
formule (ou directement de (16.7.1.2)) que
$\mathscr{P}_{X/S}^n(\mathscr{F})$ est canoniquement muni d'une structure de
$\mathscr{P}_{X/S}^n$-Module.
\end{env}

\begin{proposition}[16.7.3]
\label{IV.16.7.3}
\begin{enumerate}
\item[(i)] Le foncteur
$\mathscr{F}\mapsto\mathscr{P}_{X/S}^n(\mathscr{F})$ de la catégorie des
$\mathcal{O}_X$-Modules dans celle des $\mathscr{P}_{X/S}^n$-Modules est
exact à droite, et commute aux limites inductives quelconques; il est exact
lorsque $\mathscr{P}_{X/S}^n$ est un $\mathcal{O}_X$-Module plat.
\item[(ii)] Si $\mathscr{F}$ est un $\mathcal{O}_X$-Module quasi-cohérent
(resp. de type fini, resp. de présentation finie),
$\mathscr{P}_{X/S}^n(\mathscr{F})$ est un $\mathscr{P}_{X/S}^n$-Module
quasi-cohérent (resp. de type fini, resp. de présentation finie).
\end{enumerate}
\end{proposition}

Les assertions de (i) résultent aussitôt de la formule (16.7.2.1) et de la
considération de la symétrie de $\mathscr{P}_{X/S}^n$ (16.3.4). Les assertions
de (ii) découlent de l'exactitude à droite du foncteur
$\mathscr{F}\mapsto\mathscr{P}_{X/S}^n(\mathscr{F})$.

\begin{env}[16.7.4]
\label{IV.16.7.4}
Les deux structures de $\mathcal{O}_X$-Module de $\mathscr{P}_{X/S}^n$
définissent sur $\mathscr{P}_{X/S}^n(\mathscr{F})$ deux structures de
$\mathcal{O}_X$-Module, d'ailleurs permutables, donc une structure de
$\mathcal{O}_X$-Bimodule. Il est commode de noter \emph{à gauche} celle de
ces structures provenant de l'homomorphisme structural
$\mathcal{O}_X\to\mathscr{P}_{X/S}^n$ (choisi dans (16.3.5)) et
\emph{à droite} celle provenant de l'homomorphisme
$d_{X/S}^n:\mathcal{O}_X\to\mathscr{P}_{X/S}^n$. Autrement dit, pour tout
ouvert $U$ de $X$, et tout triplet d'éléments
$a\in\Gamma(U,\mathcal{O}_X)$, $b\in\Gamma(U,\mathscr{P}_{X/S}^n)$,
$t\in\Gamma(U,\mathscr{F})$, on a par définition
\[
  a(b\otimes t)=(ab)\otimes t,qquad
  (b\otimes t)a=(b.d^na)\otimes t=b\otimes(at)
  =(d^na).(b\otimes t).
  \tag{16.7.4.1}
  \label{IV.16.7.4.1}
\]
La structure de $\mathcal{O}_X$-Module provenant de la définition (16.7.1.2)
est donc, avec ces conventions, la structure de $\mathcal{O}_X$-Module à
gauche.

